\section{TVM：机器学习编译的开拓}

\subsection{从问题出发的动机}

从 CXXNet 到 MXNet，陈天奇发现一个反复出现的瓶颈：每一代 GPU 和新硬件都需要大量手写 kernel，支持不同数据 layout（如 NHWC vs NCHW）、不同输入尺寸、不同硬件形态，排列组合的空间巨大。手写 kernel``根本干不过来''。

\begin{importantbox}{TVM 要解决的核心问题}
能不能用编译技术自动生成这些 kernel？能不能让更少的人通过正确的技术，快速搭建面向新硬件的机器学习系统？这就是 TVM 的出发点------\textbf{让机器学习系统的工程复杂度不再随硬件种类线性增长}。
\end{importantbox}

\subsection{GPU/TPU/NPU 的区别与融合}

陈天奇在访谈中科普了三种计算芯片的区别：

\begin{itemize}
    \item \textbf{GPU}：源自图形编程，核心特性是 SIMT（Single Instruction Multiple Thread），用大量线程并行编程。随着 CUDA 的发展，已经和图形编程``打不了关系''。
    \item \textbf{TPU/NPU}：拥有更复杂的指令集，一条指令可以完成大块矩阵运算或内存读写。一般没有线程概念，因为单条指令本身就非常``大''。
    \item \textbf{融合趋势}：GPU 引入了 Tensor Core（原本是 NPU 才有的），NPU 也开始引入线程概念以接近 CUDA 生态。``现在这个界限已经越来越模糊了。''
\end{itemize}

这种融合趋势正是 TVM 持续存在价值的体现------每一代硬件的编程模型都在变化，需要编译技术来桥接差异。

\subsection{``荒岛建城''的开发过程}

TVM 的灵感部分来自 Halide（MIT 的 Jonathan Ragan-Kelley 等人的工作，在图像处理领域用编译技术优化 pipeline），但机器学习编译在当时是一片空白。陈天奇形容这段经历是``在荒岛里面建城''------没有现成的路径可以参考。

从 2016 年 4 月 GTC 回来后开始构思，到 2017 年 3 月第一个可运行版本，花了约 11 个月，基本都在写代码。他回忆这个过程是``不断迭代''的：``虽然不懂也不怕，出生牛犊不怕虎。''

在做 TVM 的同时，他和体系结构方向的合作者 Thierry Moreau 一起，用 FPGA 搭建了一个可编程的 NPU 加速器。他们最初的想法更``wild''------直接把神经网络 bake 到 FPGA 的 wire 里（后来有人做了 HLS 相关的工作）。最终他们决定先做一个类似 TPU 的 NPU，但设计理念与当时的主流完全不同。

\begin{warningbox}{ACM 班大作业经历的回报}
陈天奇在回忆 TVM 开发时多次提到 ACM 班的经历。``我们当时在 ACM 班有一个体系结构大作业，我挑战了在自己写的 CPU 上运行自己写的编译器。TVM 非常像那段经历------也是从头搭一个编译器 for 一个新的硬件。因为有了原来那段经历之后，让我觉得不是很慌。''这说明本科阶段看似``课外活动''的挑战性项目，可能在若干年后以意想不到的方式产生回报。
\end{warningbox}

陈天奇特别指出，当时他``不是做体系结构出身''，完全是 from first principle 在想怎么做。主持人评论说：``你一直都是这样------高中竞赛出生牛犊不怕虎，ACM 班写编译器出生牛犊不怕虎，做 XGBoost 出生牛犊不怕虎，做 TVM 还是出生牛犊不怕虎。这是不是你最大的优势？''陈天奇回应：``我觉得挑问题和怎么做很重要。做平庸的问题和重要的问题花的时间几乎是一样的，所以一定要挑正确的问题。''

同时，他和体系结构方向的合作者一起，用 FPGA 搭建了一个可编程的 NPU 加速器（VTA）。这个 NPU 的设计思路非常有前瞻性：

\begin{knowledgebox}{VTA 加速器的创新设计}
大部分早期 NPU 是完全 hardwired（不可编程，只有一个固定功能）。但 VTA 采用了\textbf{指令级可编程}设计，可以用同一套微指令实现多种算子。为了解决内存读写和计算的 overlap 问题，他们借鉴了超线程（hyper-threading）概念：在逻辑上将 NPU 分为两个线程，一个做内存操作，一个做计算，通过硬件 signal 机制同步。这个设计``放在今天都还挺有趣''，在 NVIDIA 最新的异步编程中能看到类似影子。
\end{knowledgebox}

\subsection{Jeff Dean 的坐标系与高风险选题}

在 Google Brain 实习期间，Jeff Dean 把陈天奇拉到会议室，画了一个坐标系：一轴是 risk，一轴是 impact。在不同的项目选项中，陈天奇选了 risk 最大、impact 也最大的课题------Net2Net。这个项目的问题是：已经有一个训练好的小模型，如何快速 bootstrap 一个大模型，使其容量可以增大？

陈天奇选择这个课题的原因是他一直有一个 vision：智能需要能演化，模型的 capacity 需要变大。一个神经网络如果一开始很小，到一定程度就会不行。从 lifelong learning 的角度，怎样解决这个问题？他指出：``其实现在你会发现这个问题还是同样存在的------这也是为什么大家开始会有通过小的模型 scaling up 的方向。''

Jeff Dean 当时是他在 Brain 的 mentor，而且``非常 open minded''。陈天奇当时算是 Jeff Dean 的第一个 intern------``他的 GAN 可能还没有发布或者刚发布''的阶段。

他的选择标准是：\textbf{``时间是有限的，所以一定要做想要做的事情。它不一定是有 impact 的事情，但一定要是想要做的事。''}

面对``你会建议年轻 researcher 选高风险高收益课题吗''的问题，他的回答很务实：不同人可能不一样。但因为他已经经历过很多失败，``再失败一次也没有那么差''。``Almost always 我们会 ends up 选一个比较高风险高回报的方向。''

\subsection{博士后期的勇气与 Recomputation 的故事}

在做 TVM 之前，陈天奇曾想到一个``用低于线性复杂度训练更深模型''的方法（通过 recomputation 节省内存），但 Carlos 说这个工作可能``发了就是个 NeurIPS poster''，不够深，要做就做到 best paper level------应该设计一个更自动化的方法来选择 recomputation 策略。因为精力有限，他没有继续这个方向，后来有 follow-up work 做了相关工作。

Carlos 的要求不是否定这个方向的价值，而是说：\textbf{``你需要挑你做的东西，并且一旦决定要做，就要做得深了。''}

李沐在知乎评论中透露：陈天奇 PhD 第三/四年跟他说要做 TVM，李沐提醒他风险太高------PhD 快结束了做一个新的大坑。陈天奇的回答是：``不管未来去工作也好、创业也好、找教职也好，最重要的还是要做当时自己觉得最重要的事。''李沐后来评价：``现在看起来这是唯一正确的做事方式。''

\begin{importantbox}{关于做选择的底层逻辑}
陈天奇做选择时会问自己一个问题：\textbf{``如果选了这条路并且失败了，我还会不会后悔？''}如果答案是不后悔，就选它。这不是理性的投入产出分析，而是确保自己追随内心。``很多时候你后悔的不是有没有获得成功，而是你是不是走了理想中想要走的那条路。''
\end{importantbox}

\subsection{TVM 的演化与生态位}

TVM 从诞生到现在已经经历了多轮重构，从最初的图编译到 Relax（为 LLM 重新设计的图编译器）。陈天奇强调 TVM 的设计哲学是\textbf{``可演化的''}------不断 reinvent 自己以适应新的需求。他指出：``你去看现在的 TVM 和它刚开始时候的形态，我们几乎已经经过了几个版本的迭代。''

面对 Torch Compile、XLA、Triton 等竞争者，他认为编译技术的需求不但没有减小，反而在增大：

\begin{itemize}
    \item \textbf{硬件 specialized}：每一代芯片往前进的方式都不同，同一个 vendor 的不同代产品编程模式都有变化
    \item \textbf{模型与硬件的``缠绕''}：为达到最优效率，模型设计越来越和硬件特性耦合------DeepSeek 针对特定硬件环境做设计就是典型例子
    \item \textbf{Hardware Lottery}：硬件限制了模型发展的可能性，需要编译技术来打破这个束缚
    \item \textbf{AI Agent 的兴起}：Agent 系统和 DSL 的出现为编译技术带来了新的应用场景
\end{itemize}

TVM 目前的发展方向有两个：一是更好地与 PyTorch 等生态整合（将通用的函数调用规范单独抽象出来，方便跨框架调用）；二是在垂直领域做突破（如将 LLM 编译部署到网页和车载设备）。他的终极目标是：\textbf{``像写神经网络一样简单地做机器学习工程开发''}------目前这还远未实现。

\begin{warningbox}{Reinvent yourself vs Patch 现有系统}
在 OctoML 内部，陈天奇也面临过 reinvent vs patch 的抉择。为了支持 LLM，他们推翻了 TVM 最初的图编译部分，重新做了 Relax。这个决定``不会被所有人认可''，因为已有技术短期可以解决当前问题。但他坚持认为：``If you don't reinvent yourself, there will be a time where you bite the technical debt。''长期来看，技术债务的积累比重构的短期成本更可怕。
\end{warningbox}

\subsection{本章小结}

TVM 是陈天奇``从问题出发''研究风格的集大成之作。它源于 MXNet 开发中发现的工程瓶颈，灵感来自编译器领域，设计中融合了硬件协同设计的思想。虽然面临激烈竞争，但编译技术在 ML 系统中的核心地位已成共识。

